%==============================================================================
% sections/10_discussion.tex
%==============================================================================
\section{Discussion and Development Roadmap}
\label{sec:discussion}

\subsection{Theoretical results established versus empirical claims projected}

The expanded manuscript separates two kinds of statements.

\paragraph{Theoretical (established under the stated assumptions).}
The independent finite-$N$ catastrophe principle
(\cref{thm:catastrophe-exact}), the single-big-jump sum equivalence
(\cref{thm:sum-equivalence}), the corrected second-order expansion with
leave-one-out mean shift (\cref{thm:second-order}), independent-CdMC
unbiasedness and the scale-invariant BRE bound $N^\alpha-1$
(\cref{prop:ind-cdmc-bre}), the arbitrary-dependence CdMC identity
(\cref{thm:dep-cdmc-unbiased}), the conditional-tail-envelope BRE result
(\cref{prop:dep-cdmc-bre}), the latent-shock tail constant and BRE bound
(\cref{thm:latent-shock-tail,prop:latent-shock-bre}), the independent-block
reduction (\cref{thm:block-reduction}), the MRV linear-risk asymptotic
(\cref{thm:mrv-linear-risk}) with the spectral CdMC identity and BRE
(\cref{prop:radial-cdmc,prop:spectral-cdmc-bre}), the hidden-cone second-order
expansion (\cref{thm:hidden-second-order}), the exact and asymptotic
control-variate VRE result (\cref{prop:vre}), and the Bernstein-type
confidence interval (\cref{thm:bernstein-ci}). Each of these is proved in
\cref{appx:independent-proofs,appx:dependent-proofs,appx:second-hidden} with
the underlying assumptions explicit.

\paragraph{Projected (planned, not yet executed).}
All simulation diagnostics in \cref{sec:simulation-plan}, all real-data
diagnostics, forecasts, and backtests in \cref{sec:real-data}, and every
figure with a (PLANNED) caption prefix or table cell marked \xx{}. Numerical
performance, backtest success, and production runtime conclusions are not
claimed until the corresponding pipeline outputs replace those entries.

\subsection{Limitations and counterexamples}
\label{subsec:limitations}

The BRE bound $N^\alpha-1$ is finite for any finite $N$ but diverges as
$N\to\infty$. The independent CdMC is therefore not suited as-is to
high-dimensional limit regimes; in such regimes the latent-shock CdMC
(\cref{prop:latent-shock-bre}) replaces $N$ by the number of active shocks
$K$, and the spectral CdMC (\cref{prop:spectral-cdmc-bre}) replaces $N$ by a
spectral integral that need not grow with the ambient dimension.

The dependent CdMC identity (\cref{thm:dep-cdmc-unbiased}) is exact only when
the conditional kernels $p_i(t;X_{-i})$ are evaluated exactly. A
misspecified copula or a poorly fit conditional kernel produces a biased
estimator at every threshold, and the bias can be much larger than the Monte
Carlo standard error. \Cref{rem:kernel-identifiability,rem:numerical-stability}
discuss identifiability and stability mitigations; the empirical protocol of
\cref{sec:real-data} reports model-residual diagnostics before any kernel
estimator is trusted.

Hidden regular variation
(\cref{def:hrv,ass:hidden-scale,thm:hidden-second-order}) requires
identification of the hidden tail scale $\Hb_2$ from finite samples. In
small samples, threshold noise can be mistaken for a hidden scale. The
manuscript treats hidden-cone results as a second-order layer and requires
threshold-stability diagnostics (\cref{appx:second-hidden}) before any
hidden-cone estimator enters a primary VaR/ES forecast.

\subsection{Connection to the literature placement}

The independent finite-$N$ sharpening of \cref{sec:independent-baseline}
positions this paper relative to \citet{PourbabaeeSolari2019} in the form
spelled out in \cref{subsec:recent-developments}: scale-invariant BRE bound,
corrected second-order expansion with $N=1$ specialization, and explicit
assumption stack. The dependent extensions of
\cref{sec:dependent-cdmcs,sec:mrv-spectral,sec:hidden-rv} connect to
\citet{AsmussenKroese2006,HartingerKortschak2009,BlanchetRojasNandayapa2011,ChengFuhPang2025}
(rare-event simulation under dependence),
\citet{HultLindskogMikoschSamorodnitsky2005,Resnick2007,SamorodnitskySun2016}
(MRV foundations), and
\citet{LedfordTawn1996,MaulikResnick2005,DasMitraResnick2013}
(hidden regular variation). The empirical protocol of \cref{sec:real-data}
follows the factor-portfolio backtesting tradition of
\citet{McNeilFreyEmbrechts2015,GlassermanHeidelbergerShahabuddin2002} and
adds dependence-diagnostic gating for estimator selection.

\subsection{Recommended development sequence}

The recommended implementation sequence is:
\begin{enumerate}
\item finalize and test the independent baseline replication;
\item implement common-shock and latent-shock simulation;
\item implement empirical marginal tail diagnostics;
\item implement pairwise tail-dependence and spectral diagnostics;
\item implement latent, block, and spectral estimators;
\item run real-data rolling VaR/ES forecasts;
\item add hidden-cone estimators only after diagnostics justify them.
\end{enumerate}

\subsection{Conclusion}

The original independent-factor paper is strongest when the extreme event is
well represented by one large observed coordinate. The financial problem asks
for more: extremes may be generated by one latent market shock, one dependent
block, one radial vector movement, or one hidden joint-tail cone. The
expanded paper keeps the original single-big-jump insight but moves it to
the correct object. In the independent model that object is a coordinate;
in the latent model it is a shock; in MRV it is a spectral/radial movement;
and in hidden regular variation it is a slower-scale joint cone. The
accompanying empirical protocol is designed to determine which object is
visible in real financial data and whether modeling it improves VaR, ES,
attribution, and simulation efficiency. The full method is implemented as
the open-source \texttt{FactorTail} library
(\url{https://github.com/osolari/FactorTail}); the repository is the
canonical reference for any future numerical or empirical claim downstream
of this manuscript.
